%=======================================================================
% CONTINUAÇÃO DO TCC_V2.tex - PARTE 2
% Adicionar esta seção ao final do arquivo tcc_v2.tex
%=======================================================================

%=======================================================================
% 6 LIMITAÇÕES E TRABALHOS FUTUROS
%=======================================================================
\chapter{Limitações e Trabalhos Futuros}

\section{Limitações Identificadas}

\subsection{Limitações Técnicas}
Durante o desenvolvimento e teste do sistema, foram identificadas as seguintes limitações:

\begin{itemize}
    \item \textbf{Dependência de Conectividade:} O sistema requer conexão Wi-Fi estável para funcionalidades remotas, embora mantenha controle local básico durante intermitências.
    \item \textbf{Calibração Específica:} Sensores capacitivos necessitam calibração individual por tipo de solo e vaso, limitando a portabilidade da configuração.
    \item \textbf{Escalabilidade:} A arquitetura atual suporta um dispositivo por instância Firebase, requerendo modificações para múltiplos dispositivos.
    \item \textbf{Precisão Climática:} Sensores de baixo custo (DHT22, LDR) apresentam limitações de precisão em condições extremas.
    \item \textbf{Capacidade de Bomba:} Sistema atual adequado para vasos pequenos/médios (até 5L), necessitando bombas mais potentes para aplicações maiores.
\end{itemize}

\subsection{Limitações de Usabilidade}
\begin{itemize}
    \item \textbf{Interface Web Única:} Não possui aplicativo móvel nativo, dependendo do navegador web.
    \item \textbf{Configuração Manual:} Requer conhecimento básico sobre necessidades hídricas das plantas para configuração ótima.
    \item \textbf{Manutenção:} Necessita limpeza periódica dos sensores e verificação do reservatório de água.
\end{itemize}

\subsection{Limitações Experimentais}
\begin{itemize}
    \item \textbf{Período de Teste:} Experimentos limitados a 7-14 dias não capturam variações sazonais.
    \item \textbf{Diversidade de Plantas:} Testes focados em espécies ornamentais comuns, não incluindo plantas com necessidades hídricas específicas.
    \item \textbf{Ambiente Controlado:} Testes realizados em ambiente interno podem não refletir condições externas variáveis.
\end{itemize}

\section{Trabalhos Futuros}

\subsection{Melhorias de Hardware}
\begin{itemize}
    \item \textbf{Conectividade Redundante:} Implementação de conectividade LoRa ou 4G para backup do Wi-Fi doméstico.
    \item \textbf{Sensores Avançados:} Integração de sensores de condutividade (EC) e pH para monitoramento nutricional completo.
    \item \textbf{Alimentação Solar:} Sistema de energia renovável com bateria para operação autônoma.
    \item \textbf{Múltiplos Vasos:} Expansão para controle simultâneo de 4-8 vasos com válvulas individuais.
    \item \textbf{Sensores Ambientais:} Adição de barômetro e anemômetro para correlação com dados meteorológicos.
\end{itemize}

\subsection{Aprimoramentos de Software}
\begin{itemize}
    \item \textbf{Inteligência Artificial:} Implementação de modelos de aprendizado de máquina para predição de necessidades hídricas baseada em dados históricos, clima e tipo de planta.
    \item \textbf{Integração Meteorológica:} Consumo de APIs climáticas (OpenWeatherMap) para ajuste automático baseado em previsão de chuva e temperatura.
    \item \textbf{Aplicativo Móvel:} Desenvolvimento de app nativo Flutter para iOS/Android com notificações push.
    \item \textbf{Sistema de Alertas:} Integração com WhatsApp, Telegram ou email para notificações de manutenção e alertas críticos.
    \item \textbf{Analytics Avançado:} Dashboard com análise de tendências, economia de água acumulada e recomendações personalizadas.
\end{itemize}

\subsection{Funcionalidades Avançadas}
\begin{itemize}
    \item \textbf{Reconhecimento de Plantas:} Uso de visão computacional para identificação automática de espécies e configuração de parâmetros.
    \item \textbf{Sistema de Recompensas:} Gamificação para incentivar práticas sustentáveis de jardinagem.
    \item \textbf{Marketplace de Configurações:} Plataforma para compartilhamento de perfis de plantas entre usuários.
    \item \textbf{Integração Home Assistant:} Compatibilidade com sistemas de automação residencial existentes.
    \item \textbf{Controle por Voz:} Integração com Alexa/Google Assistant para comandos de voz.
\end{itemize}

\subsection{Pesquisa e Desenvolvimento}
\begin{itemize}
    \item \textbf{Algoritmos Adaptativos:} Desenvolvimento de algoritmos que se ajustam automaticamente às características específicas de cada planta e ambiente.
    \item \textbf{Estudo Longitudinal:} Experimentos de longo prazo (6-12 meses) para validação em diferentes estações.
    \item \textbf{Eficiência Energética:} Otimização do consumo de energia através de técnicas de deep sleep e harvesting de energia.
    \item \textbf{Padrões IoT:} Implementação de protocolos padronizados (MQTT, CoAP) para interoperabilidade.
    \item \textbf{Edge Computing:} Processamento local de dados para reduzir dependência de conectividade.
\end{itemize}

\section{Impacto e Escalabilidade}

\subsection{Potencial de Mercado}
O sistema desenvolvido possui potencial de aplicação em diversos segmentos:
\begin{itemize}
    \item \textbf{Residencial:} Jardins domésticos, varandas e hortas urbanas
    \item \textbf{Educacional:} Projetos escolares de sustentabilidade e tecnologia
    \item \textbf{Comercial:} Pequenas estufas e viveiros
    \item \textbf{Social:} Projetos comunitários de agricultura urbana
\end{itemize}

\subsection{Replicabilidade}
A documentação completa e uso de componentes padronizados facilitam a replicação do sistema por:
\begin{itemize}
    \item Estudantes de engenharia e ciência da computação
    \item Entusiastas de IoT e automação residencial
    \item Produtores rurais de pequena escala
    \item Projetos de extensão universitária
\end{itemize}

\subsection{Contribuição Científica}
Este trabalho contribui para a área de IoT aplicada à agricultura através de:
\begin{itemize}
    \item Metodologia de calibração para sensores capacitivos de umidade
    \item Arquitetura de referência para sistemas de irrigação doméstica
    \item Protocolo experimental para avaliação de sistemas IoT de irrigação
    \item Análise comparativa de tecnologias de conectividade para IoT agrícola
\end{itemize}

%=======================================================================
% 7 ASPECTOS AMBIENTAIS E SUSTENTABILIDADE
%=======================================================================
\chapter{Aspectos Ambientais e Sustentabilidade}

\section{Impacto Ambiental}

\subsection{Conservação de Recursos Hídricos}
O sistema proposto contribui para a sustentabilidade através de:
\begin{itemize}
    \item \textbf{Redução do Desperdício:} Economia média de 33\% no consumo de água comparado à irrigação manual tradicional.
    \item \textbf{Irrigação Precisão:} Fornecimento de água apenas quando necessário, evitando encharcamento e runoff.
    \item \textbf{Otimização Temporal:} Irrigação em horários apropriados para minimizar evaporação.
    \item \textbf{Monitoramento Contínuo:} Dados históricos permitem identificar padrões e otimizar consumo.
\end{itemize}

\subsection{Eficiência Energética}
\begin{itemize}
    \item \textbf{Baixo Consumo:} ESP32 consome ~240mA em operação ativa, <10µA em deep sleep.
    \item \textbf{Componentes Eficientes:} Bomba de baixa tensão (3-6V) reduz consumo energético.
    \item \textbf{Controle Local:} Reduz dependência de comunicação constante com a nuvem.
    \item \textbf{Potencial Solar:} Arquitetura compatível com alimentação por painéis solares pequenos.
\end{itemize}

\subsection{Redução de Insumos Químicos}
\begin{itemize}
    \item \textbf{Saúde das Plantas:} Irrigação adequada reduz stress e necessidade de defensivos.
    \item \textbf{Prevenção de Doenças:} Evita encharcamento que favorece fungos e pragas.
    \item \textbf{Crescimento Otimizado:} Plantas saudáveis requerem menos fertilizantes corretivos.
\end{itemize}

\section{Análise de Ciclo de Vida}

\subsection{Produção e Materiais}
\begin{itemize}
    \item \textbf{Componentes Duráveis:} Sensores capacitivos resistem à corrosão, aumentando vida útil.
    \item \textbf{Materiais Recicláveis:} PCBs e componentes eletrônicos são recicláveis através de programas específicos.
    \item \textbf{Embalagem Mínima:} Componentes individuais reduzem desperdício de embalagem.
\end{itemize}

\subsection{Operação}
\begin{itemize}
    \item \textbf{Vida Útil Estimada:} 5-7 anos para componentes eletrônicos principais.
    \item \textbf{Manutenção Mínima:} Limpeza periódica de sensores e verificação de conexões.
    \item \textbf{Atualizações OTA:} Firmware atualizável remotamente prolonga vida útil.
\end{itemize}

\subsection{Descarte}
\begin{itemize}
    \item \textbf{Reciclagem Eletrônica:} Componentes podem ser encaminhados para reciclagem apropriada.
    \item \textbf{Reutilização:} Sensores e microcontrolador podem ser reutilizados em outros projetos.
    \item \textbf{Biodegradabilidade:} Componentes não eletrônicos (mangueiras, conectores) podem ser substituídos por materiais biodegradáveis.
\end{itemize}

\section{Educação e Conscientização}

\subsection{Valor Educativo}
O sistema serve como ferramenta educativa para:
\begin{itemize}
    \item \textbf{Conceitos de IoT:} Demonstração prática de conectividade e automação.
    \item \textbf{Sustentabilidade:} Visualização concreta de economia de recursos.
    \item \textbf{Agricultura 4.0:} Introdução a conceitos de agricultura de precisão.
    \item \textbf{Programação:} Código aberto permite aprendizado de desenvolvimento embarcado.
\end{itemize}

\subsection{Engajamento Comunitário}
\begin{itemize}
    \item \textbf{Projetos Escolares:} Adaptável para feiras de ciências e projetos interdisciplinares.
    \item \textbf{Workshops:} Material para oficinas de tecnologia e sustentabilidade.
    \item \textbf{Replicação:} Documentação completa facilita implementação por terceiros.
    \item \textbf{Dados Abertos:} Possibilidade de contribuir para pesquisas sobre irrigação urbana.
\end{itemize}

%=======================================================================
% 8 CONCLUSÃO
%=======================================================================
\chapter{Conclusão}

Este trabalho apresentou o desenvolvimento e avaliação de um sistema de irrigação automatizado para plantas domésticas baseado em tecnologias de Internet das Coisas. A solução integra sensores de umidade do solo, temperatura, umidade do ar e luminosidade a um microcontrolador ESP32, utilizando Firebase Realtime Database como backend e uma interface web responsiva para monitoramento e controle remoto.

\section{Objetivos Alcançados}

Os objetivos propostos foram integralmente atendidos:

\textbf{Monitoramento em Tempo Real:} O sistema realiza monitoramento contínuo da umidade do solo e variáveis ambientais (temperatura, umidade do ar, luminosidade) com atualizações em tempo real via Firebase Realtime Database.

\textbf{Controle Automático Adaptativo:} Implementou-se algoritmo de controle baseado em máquina de estados com histerese, permitindo configuração de parâmetros específicos por espécie de planta através de perfis pré-definidos ou configuração manual.

\textbf{Interface Intuitiva:} A interface web desenvolvida oferece visualizações gráficas em tempo real, histórico de dados via Chart.js, gerenciamento de perfis de plantas e comandos remotos, com design responsivo para diferentes dispositivos.

\textbf{Configuração Flexível:} O sistema permite tanto configuração simples através de sliders quanto configuração avançada com parâmetros técnicos detalhados, atendendo usuários com diferentes níveis de conhecimento técnico.

\section{Contribuições}

As principais contribuições deste trabalho incluem:

\textbf{Arquitetura Integrada:} Desenvolvimento de arquitetura IoT completa integrando ESP32, sensores múltiplos, Firebase e interface web com comunicação bidirecional em tempo real.

\textbf{Controle de Potência Robusto:} Implementação de circuito de acionamento usando transistor TIP127 (PNP) para controle confiável de bomba d'água com proteções de segurança.

\textbf{Sistema de Calibração:} Metodologia de calibração para sensores capacitivos de umidade usando interpolação linear com filtragem por média móvel.

\textbf{Interface Adaptativa:} Desenvolvimento de interface web com dois modos de operação (simples/avançado) e funcionalidades avançadas de monitoramento climático.

\textbf{Protocolo Experimental:} Estabelecimento de metodologia reproduzível para avaliação de sistemas de irrigação IoT com métricas quantitativas de economia de água e estabilidade.

\textbf{Documentação Completa:} Produção de documentação técnica detalhada incluindo esquemas elétricos, código fonte e instruções de montagem para facilitar replicação.

\section{Resultados Significativos}

Os resultados experimentais demonstraram a eficácia da solução:

\begin{itemize}
    \item \textbf{Economia de Água:} 33.3\% de redução no consumo comparado à irrigação manual
    \item \textbf{Estabilidade da Umidade:} Redução do desvio padrão de 8.7\% para 3.1\%
    \item \textbf{Tempo na Faixa Alvo:} Melhoria de 56\% para 89\% do tempo dentro da faixa ótima
    \item \textbf{Responsividade:} Latência média de 150ms para comandos remotos
    \item \textbf{Disponibilidade:} 98.5\% de uptime durante período de teste
    \item \textbf{Custo Competitivo:} R\$ 140,00 comparado a R\$ 300-800 de soluções comerciais
\end{itemize}

\section{Impacto e Relevância}

O sistema desenvolvido apresenta relevância em múltiplas dimensões:

\textbf{Sustentabilidade:} Contribui para uso racional de recursos hídricos e redução de desperdício, alinhado com objetivos de desenvolvimento sustentável.

\textbf{Acessibilidade:} Custo reduzido e componentes amplamente disponíveis tornam a tecnologia acessível para uso doméstico.

\textbf{Educacional:} Serve como plataforma educativa para conceitos de IoT, programação embarcada e agricultura de precisão.

\textbf{Escalabilidade:} Arquitetura modular permite expansão para múltiplos vasos e funcionalidades adicionais.

\textbf{Replicabilidade:} Documentação completa e código aberto facilitam adoção e adaptação por terceiros.

\section{Perspectivas Futuras}

O trabalho estabelece base sólida para desenvolvimentos futuros significativos:

\textbf{Inteligência Artificial:} Integração de modelos de aprendizado de máquina para predição adaptativa de necessidades hídricas.

\textbf{Conectividade Avançada:} Implementação de protocolos IoT padronizados (MQTT, LoRaWAN) para maior interoperabilidade.

\textbf{Sensoriamento Ampliado:} Adição de sensores de pH, condutividade e macro/micronutrientes para monitoramento nutricional completo.

\textbf{Integração Meteorológica:} Consumo de dados climáticos para ajuste automático baseado em previsões meteorológicas.

\textbf{Aplicativo Móvel:} Desenvolvimento de aplicação nativa para iOS/Android com notificações push e funcionalidades offline.

\section{Considerações Finais}

Este trabalho demonstrou a viabilidade técnica e econômica de sistemas de irrigação automatizada baseados em IoT para aplicações domésticas. A integração de tecnologias acessíveis (ESP32, Firebase, sensores de baixo custo) resultou em solução funcional que atende às necessidades práticas de usuários domésticos enquanto contribui para práticas sustentáveis de jardinagem.

A metodologia desenvolvida e os resultados obtidos estabelecem fundamentos para trabalhos futuros em agricultura urbana de precisão e IoT aplicada à sustentabilidade. A disponibilização do código fonte e documentação técnica visa facilitar a replicação e evolução da solução por pesquisadores e entusiastas da área.

O sistema proposto representa passo significativo na democratização de tecnologias de irrigação inteligente, tornando-as acessíveis para uso residencial e contribuindo para práticas mais sustentáveis de cultivo doméstico. Os resultados positivos de economia de água e melhoria na estabilidade da umidade do solo validam a abordagem proposta e incentivam desenvolvimentos futuros na área.

%=======================================================================
% REFERÊNCIAS BIBLIOGRÁFICAS
%=======================================================================
\chapter{Referências}

% Nota: Estas referências são exemplares para demonstrar a formatação.
% Em um trabalho real, seria necessário incluir as referências efetivamente utilizadas.

\bibliographystyle{abntex2-alf}
\begin{thebibliography}{99}

\bibitem{silva2019}
SILVA, João A.; SANTOS, Maria B. \textbf{Internet das Coisas aplicada à automação residencial}: fundamentos e aplicações práticas. São Paulo: Editora Técnica, 2019. 245p.

\bibitem{fernandes2021}
FERNANDES, Carlos R.; OLIVEIRA, Ana P.; COSTA, Pedro L. Sistemas de irrigação automatizada para agricultura de precisão: revisão sistemática e tendências. \textbf{Revista Brasileira de Engenharia Agrícola e Ambiental}, v. 25, n. 3, p. 187-195, 2021.

\bibitem{pereira2021}
PEREIRA, Lucas M.; RODRIGUES, Fernanda S. Eficiência hídrica em sistemas de irrigação inteligente: estudo comparativo. \textbf{Engenharia Agrícola}, v. 41, n. 2, p. 223-234, 2021.

\bibitem{santos2020}
SANTOS, Rafael T.; LIMA, Juliana C.; SOUSA, Marcos V. ESP32 em aplicações IoT: características técnicas e implementação prática. \textbf{IEEE Latin America Transactions}, v. 18, n. 5, p. 892-901, 2020.

\bibitem{almeida2020}
ALMEIDA, Gustavo H.; BARBOSA, Letícia F. Sensores capacitivos para monitoramento de umidade do solo: calibração e validação. \textbf{Sensors and Actuators A: Physical}, v. 315, p. 112-119, 2020.

\bibitem{firebase2023}
GOOGLE LLC. \textbf{Firebase Realtime Database Documentation}. Disponível em: https://firebase.google.com/docs/database. Acesso em: 15 out. 2024.

\bibitem{volkmann2022}
VOLKMANN, Andreas; MÜLLER, Klaus; WAGNER, Sabine. Low-cost IoT irrigation system with MQTT communication and mobile interface. \textbf{Computers and Electronics in Agriculture}, v. 195, p. 106-118, 2022.

\bibitem{medeiros2021}
MEDEIROS, Paulo R.; NASCIMENTO, Carla M.; XAVIER, Roberto A. Sistema de irrigação automatizada para plantas domésticas baseado em Arduino e aplicação web. \textbf{Revista de Engenharia e Pesquisa Aplicada}, v. 6, n. 2, p. 45-56, 2021.

\bibitem{lima2019}
LIMA, Sandra J.; MOREIRA, Antônio C.; SILVA, Eduardo P. Automação da irrigação na agricultura familiar: revisão da literatura e perspectivas futuras. \textbf{Revista Brasileira de Agricultura Irrigada}, v. 13, n. 4, p. 3542-3558, 2019.

\end{thebibliography}

\end{document}
